\subsection*{受限接口为何有效：表达力与性能}

逐条编写装入和地址计算给程序员很多选择，也让编译器需要重新识别哪些操作共同构成一次分块搬运。Croqtile 要求程序显式描述这片区域。当所需搬运能够用规则形状表示时，这项约束很有用，因为在生成各次传输之前，区域几何信息就已明确。

Croqtile 用结构化搬运操作显式描述分块形状和布局。编译器据此生成传输代码，程序员也不必为每个元素反复编写地址计算与协作逻辑。

图~\ref{fig:f09-fig8_expressiveness} 对比了「能写出来的范围」与「跑得快的范围」。

\begin{figure}[htbp]\centering
\includegraphics[width=\linewidth,height=0.56\textheight,keepaspectratio]{f09-fig8_expressiveness_light.png}
\caption{结构化分块操作显式提供形状和布局，便于编译器生成高效代码；实际性能仍取决于分块选择与执行配置。}
\label{fig:f09-fig8_expressiveness}
\end{figure}

这里的取舍是用明确的分块表达替代任意的逐元素搬运代码，从而便于高效生成目标代码。实际性能仍取决于分块尺寸、对齐方式、布局和同步安排，并不由语法单独保证。
